% SPDX-License-Identifier: CC-BY-SA-4.0
% Author: Matthieu Perrin

\begingroup

\SetKwFunction{Propose}{propose}
\SetKwData{BinCons}{bin\_cons}
\SetKwVariable{Proposals}{proposals}
\SetKwMessage{PROP}{proposal}
\SetKwBroadcast{URB}{urb}
\SetKwBroadcast{RB}{rb}
 
\begin{exercice}[Du consensus binaire au consensus multi-valué\footnote{[MRT00] Achour Mostefaoui, Michel Raynal and Frédéric Tronel. \emph{From binary consensus to multivalued consensus in asynchronous message-passing systems.} IPL 2000}]
 
  Le but de cet exercice est d'implémenter le consensus multi-valué à partir du consensus binaire. 
  \begin{description}
  \item[Consensus binaire :] les valeurs proposées et décidées appartiennent à $\{\True,\False\}$.
  \item[Consensus multi-valué :] les valeurs proposées et décidées appartiennent à un ensemble de valeurs arbitraire.
  \end{description}
  
  On considère l'algorithme~\ref{algo:binConsTomvCons} défini ci-dessous.
  
  \begin{algorithm}
    \lSVariables{}{
      $\BinCons[0, 1, 2, ...]$\tcp*[f]{Séquence infinie d'instances du consensus binaire}
    }
    \lLVariables{\At $p_i$}{
      $\Proposals_i \leftarrow [\bot, ..., \bot]$\tcp*[f]{$\Proposals_i[j] = $ la proposition de $p_j$, si $p_i$ la connaît}
    }
    \Operation{$\Propose(v)$ \At $p_i$}{
      \nl \URB.\Broadcast $\PROP(v)$\;
      \nl \For{$r = 0, 1, 2, ...$\label{algo:binConsTomvCons:for}}{
        \nl \Let $j = (r \mod n) + 1$\;
        \nl \If{$\BinCons[r].\Propose(\Proposals_i[j] \neq \bot)$}{
          \nl \Wait{$\Proposals_i[j] \neq \bot$}\label{algo:binConsTomvCons:wait}\;
          \nl \Return $\Proposals_i[j]$\;
        }
      }
    }
    \nl \lWhen{\URB.\Deliver $\PROP(v)$ \From $p_j$ \At $p_i$}{
      $\Proposals_i[j] \leftarrow v$
    }
    \caption{Du consensus binaire au consensus multi-valué}
    \label{algo:binConsTomvCons}
  \end{algorithm}


  \begin{questions}

  \item On considère une ronde $r$ de la boucle pour (ligne~\ref{algo:binConsTomvCons:for}),
    et la valeur $j=(r\bmod n)+1$ correspondant, pour un processus $p_i$.
    Expliquer ce que signifie proposer respectivement $\True$ ou $\False$ à l'instance $\BinCons[r]$
    de consensus binaire. Que signifie une décision $\True$ ?

  \item Donner une exécution dans laquelle un processus attend à la ligne~\ref{algo:binConsTomvCons:wait}.
    Expliquer pourquoi cette attente est nécessaire.

  \item Pourquoi l'algorithme considère-t-il une infinité d'instances de consensus
    binaire en parcourant cycliquement les $n$ processus ?
    Montrer qu'une variante qui ne considérerait chaque processus qu'une seule
    fois pourrait ne jamais décider.

  \item On propose de remplacer la diffusion uniforme fiable \URB-broadcast des
    propositions par une diffusion fiable \RB-broadcast.
    Construire une exécution dans laquelle un processus correct ne termine
    jamais son appel à $\Propose$.
    
  \item En déduire que tous les processus qui décident retournent la même valeur,
    que cette valeur a été proposée, et que tout processus correct finit par
    décider.

  \end{questions}
  
\end{exercice}

\endgroup
\endinput
